\chapter{Introduction}

\emph{Feature selection} is an important task in the machine learning and statistics community. In standard supervised classification or regression problems, we are given a number observations with, usually, multivariate features (independent variables) and a target variable or class as the dependent variable. The learning task comprises of fitting a mathematical model based on this training data set which is then used for prediction of the target variable for unseen data points. The key measure for this task is to usually reduce the error in prediction or penalty for misprediction. But it is often observed that many of such input features do not play any significant role in the mathematical model being learned, either by not being sufficiently correlated with the target variable and therefore being useless or by being highly correlated with another feature and therefore being redundant. In such scenarios, feature selection task can be used as a pre-processing step for the data set, in which we select, based on some measure, the useful features and get rid of the rest before the learning task begins. Therefore, from system's perspective, essentially feature selection pre-processing can dramatically reduce the complexity of the learning problem. It also helps inferring the model manually where understanding the roles of the input features to predict the target in the model is desired.

\section{Feature selection - overview}
Selecting features from a given set based on certain measures is a combinatorial optimization problem and is usually hard. A naive approach towards feature selection can be a cross validation setting, where all possible combinations of features are used to learn models for prediction and the features corresponding to the model with minimal prediction error are selected. However such a system is infeasible due to the exponential nature of this approach. Also, it requires invoking the learner itself in the pre-processing step and therefore defeats the purpose of feature selection in the first place. Usual methods thereby aim for minimizing a ``proxy measure'' for prediction error. The two most commonly used algorithms for this purpose are \emph{Forward Selection} and \emph{Backward Elimination}. In Forward Selection case, the approach starts with an empty set in hand and it iteratively keeps adding new features to that set minimizing (or maximizing) some measure until the desired number of features is reached (which usually is a user input). Backward Elimination, on the other hard, takes the same approach from the other direction, \textit{i.e.} start with the entire set and iteratively gets rid of the features one by one until desired number of features are left in the set.

\subsection{Heuristic for feature selection}
Finding the proper measure for feature selection algorithms is fundamentally a heuristic and is specific to the mathematical nature of the model since the notion of dependency or correlation depends on that. For example, consider that a support vector machine (SVM) classifier with some string kernel is being used as the model in order to classify e-mails as spam or ham. A feature selection task here can be to induce a set of keywords which play important role in such classification. Therefore, the appropriate measure for this task should be to look for dependency and correlation in the kernel mapped space. Linear regression, on the other hand, is one of the simplest and widely used methods where the mathematical model used for fitting is linear in the features and the target variable. The notion of dependency between the variables is easily captured by the covariance or correlation matrix which is easy to compute. Therefore, a popular measure for feature selection for linear regression is usually the squared multiple correlation ($R^2$-statistic) which can be computed in terms of the covariance matrices when the data is normalized (see \citealp{Das:2008:ASS:1374376.1374384}). Another widely used algorithm for feature selection in linear settings is Least Absolute Shrinkage and Selection Operator (LASSO) which uses $l_1$-norm based regularizer to promote sparsity in the selected features.

\subsection{Diversity in the selected features}
Apart from the measure that focuses on the dependency of the selected features and the target variable, another important aspect of feature selection is to capture diversity in the selected features in order to minimize redundancy. Diverse features often make more sense from interpreter's point of view as well. In addition, diversity in features often makes the model robust in presence of errors in the training data. Therefore, usage of diversity promoting regularizers in feature selection task is novel and it performs extremely well in many practical problems (see \citealp{NIPS2012_4689}).

\subsection{Usage of submodularity in feature selection}
Submodularity is an elegant concept in mathematics that is relatively new but it enjoys the importance from many diverse fields that root in mathematical foundation. In \citealp{liu-submodfeature2013-icassp}, the concept of submodularity was used in a feature selection algorithm that tries to select a subset of features that contains as much information as possible while being much smaller than the original set. Earlier in \citealp{DasDK11} in the context of linear regression, it is shown that in feature selection task, while achieving exact submodularity using $R^2$ objective function is not possible, we can still get some optimality guarantee using the notion of \emph{approximate submodularity}. Then in \citealp{NIPS2012_4689}, a number of \emph{submodular regularizers} were used along with the same objective function to promote diversity in the selected features. In that work, they had shown that while using such regularizers, we can still achieve some optimality guarantee for that objective function. 

As for monotone submodular regularizers, \citealp{NIPS2012_4689} had used a simple Forward Selection algorithm, but the non-monotone case was rather complex. Since the goal for this particular setting is to maximize the objective function, the presence of non-monotone submodular regularizers makes the task hard. But in that work they had shown that using a local search algorithm similar to \citealp{Feige:2011:MNS:2079073.2079078} along with Forward Selection we can still reach an optimality guarantee. Even though the local search algorithm was originally designed for unconstrained maximization of non-monotone \emph{submodular} objective, using a slight variation of that algorithm helped proving an optimality guarantee for their original objective function as well which is \emph{not submodular} but \emph{has a submodular regularizer} component. While the performance of this algorithm is extremely well, the complexity of this algorithm is relatively high due to the expensive local search algorithm. 

In this present work, we explore the usage of a \emph{linear time} algorithm for local search based on \citealp{Buchbinder:2012:TLT:2417500.2417835} which was designed for unconstrained submodular maximization (USM). It is interesting to note that, for using such an algorithm, the key idea is to \emph{draw a relation} between \emph{the solution set} returned by the original optimization algorithm (which works on a \emph{submodular} objective) and \emph{the optimal set} for the final objective function (whose optimal set can be different than the optimal set of the submodular component). For the local search based on \citealp{Feige:2011:MNS:2079073.2079078}, this connection was explored in detail in the original paper. The main challenge in our approach was to obtain such a connection while using the linear time variant, which was not obvious from the description of the algorithm. The main contribution of this work is theorem \ref{thm:mine}, which forms the base of using  a linear time algorithm for a generic objective function with submodular component. Following that, we show that we can still reach a provable optimality guarantee using this linear time variant, therefore reduce the overall complexity of the whole algorithm for diverse feature selection problem in linear regression setting.

\section{Structure of the report}
The content of this report is organized as follows. In chapter \ref{chapter:background}, we review the key mathematical concepts regarding submodularity, approximate submodularity and submodularity ratio. We also discuss the feature selection task for linear regression problem in detail. We review the squared multiple correlation or $R^2$-objective function and diversity promoting regularizers introduced in \citealp{NIPS2012_4689}. In chapter \ref{chapter:work}, we discuss the local search algorithms for USM and their complexity and optimality guarantee. We discuss our theoretical results in this chapter as well present with proofs the optimality guarantee of our approach. Then in chapter \ref{chapter:results}, we discuss our experiments and results. Finally in chapter \ref{chapter:remarks}, we put concluding remarks for this work.
